% New research fragment. Root and path_stability independently audited the safe exponential bound.
% Do not merge into the previously audited manuscripts without a separate review.
\section{Exponential decay for disjoint regions with bounded cuts}
\label{sec:second-order-exponential-cuts}

All cycle spaces below are over $\mathbb F_2$.
Graphs may have parallel edges.
For a graph $Q$, write $\mathcal Z(Q)$ for its binary cycle space.
A uniform element of this space is denoted by $F$.
A vertex is \emph{closed} when no incident edge belongs to $F$.
Loops can be deleted when discussing this event, since it concerns the edge boundary of a vertex.

\begin{lemma}[A common local functional]\label{lem:common-star-functional}
Let $Q$ be a connected loopless multigraph.
Let $v_1,\ldots,v_\ell$ be distinct vertices of degree at most $d$, where $d\ge2$.
Suppose a nonzero functional $a\in\mathcal Z(Q)^*$ lies in the span of the incident edge-coordinate functionals at every $v_i$.
If $Z$ counts the closed vertices among $v_1,\ldots,v_\ell$, then
\[
 \mathbb E\,2^{-Z}
 \le\frac12+\frac12\left(1-2^{1-d}\right)^{\lfloor\ell/2\rfloor}.
\]
\end{lemma}

\begin{proof}
At $v_i$, choose an incident-edge subset $A_i$ whose coordinate sum represents $a$.
It is nonempty and proper: the empty subset and the whole star both represent zero.
Split $v_i$ into two vertices carrying $A_i$ and its complement, and join them by a new edge $h_i$.
Perform all splits simultaneously, obtaining a connected graph $Q'$.
Contracting the new edges identifies the two cycle spaces.
Under this identification, the coordinate on every $h_i$ equals $a$.

Deleting all $\ell$ new edges imposes exactly the one condition $a=0$.
Thus
\[
 \beta\bigl(Q'-\{h_1,\ldots,h_\ell\}\bigr)=\beta(Q')-1.
\]
Keeping isolated vertices in the rank calculation shows that the remaining graph has exactly $\ell$ components.
Contract these components.
The resulting connected quotient has $\ell$ vertices and $\ell$ edges.
An even subgraph with $a=1$ projects to its entire edge set.
The quotient is therefore connected and two-regular, hence a cycle.
For $\ell=1$ it is a loop, and for $\ell=2$ it consists of two parallel edges.

Condition on $a=0$.
The restrictions to the $\ell$ components are independent uniform even subgraphs.
Choose a matching of $\lfloor\ell/2\rfloor$ edges in the quotient cycle.
For each chosen $h_i$, its two split endpoints belong to distinct components.
Different chosen edges use disjoint component pairs.

Suppose the two parts of the original star have sizes $s_i$ and $t_i$.
Their sum is $\deg_Q(v_i)\le d$.
In its component, a vertex of degree $s_i$ has at most $s_i-1$ independent incident-edge coordinates.
Its star is therefore zero with probability at least $2^{1-s_i}$.
The corresponding bound at the other endpoint is $2^{1-t_i}$.
Consequently, the original vertex $v_i$ is closed with conditional probability at least $2^{2-d}$.
For the chosen matching, these events are independent.
It follows that
\[
 \mathbb E[2^{-Z}\mid a=0]
 \le\left(1-2^{1-d}\right)^{\lfloor\ell/2\rfloor}.
\]
When $a=1$, none of the selected original vertices is closed.
Since $a$ is a nonzero functional, its two values have equal probability.
This proves the result.
For $\ell=1$, the displayed bound is the trivial bound one.
\end{proof}

\begin{theorem}[Exponential closed-star estimate]\label{thm:exponential-closed-stars}
Let $Q$ be a connected loopless multigraph with $k\ge1$ distinct marked vertices of degree at most $d$.
Let $Z$ count the marked vertices closed in a uniform even subgraph.
For $d\ge2$, put
\[
 T_d=2^d,\qquad D_d=(d-1)(2^d-1),\qquad M_d=2^{D_d}-1.
\]
Then
\begin{equation}\label{eq:exponential-closed-stars}
 \mathbb E\,2^{-Z}
 \le\frac12+\exp\!\left(-\frac{k}{2^d M_d}\right)
 \le\frac12+\exp\!\left(-\frac{k}{2^{D_d+d}}\right).
\end{equation}
For $d\le1$, every marked vertex is closed and $\mathbb E2^{-Z}=2^{-k}$.
\end{theorem}

\begin{proof}
Let $L_i\subseteq\mathcal Z(Q)^*$ be the span of the edge-coordinate functionals at marked vertex $i$.
The parity equation at that vertex gives $\dim L_i\le d-1$.
If some $L_i$ is zero, that vertex is always closed.
The expectation is then at most one half, so assume every $L_i$ is nonzero.

If $T_d$ of these spaces form a direct sum, their local patterns are independent.
Each corresponding vertex is closed with probability at least $2^{1-d}$.
Consequently,
\[
 \mathbb E2^{-Z}\le(1-2^{-d})^{T_d}\le e^{-1}<\frac12.
\]

Otherwise choose a maximal direct-sum family and let $U$ be its span.
It has fewer than $T_d$ members, so $\dim U\le D_d$.
Every unselected $L_i$ meets $U$ nontrivially, by maximality.
Every selected $L_i$ also meets $U$ nontrivially, since it is nonzero and contained in $U$.
There are at most $M_d$ nonzero vectors in $U$.
Thus one nonzero functional $a$ belongs to at least $\ell\ge k/M_d$ of the spaces $L_i$.

Apply Lemma~\ref{lem:common-star-functional} to those vertices and ignore the other marked factors.
With $q=1-2^{1-d}\ge1/2$, one has
\[
 \frac12q^{\lfloor\ell/2\rfloor}
 \le\frac1{2q}\exp\!\left(-\frac{(1-q)\ell}{2}\right)
 \le\exp(-2^{-d}\ell).
\]
This proves~\eqref{eq:exponential-closed-stars}.
The assertion for $d\le1$ follows directly from parity.
\end{proof}

\begin{corollary}[Disjoint cyclic regions]\label{cor:exponential-small-cuts}
Let $G$ be a connected graph and let $U_1,\ldots,U_k$ be pairwise disjoint vertex sets.
Suppose each $G[U_i]$ is connected and contains a cycle, and $|\delta_G(U_i)|\le d$.
For a uniform even subgraph $F$ of $G$, let $\mathcal F$ be the event that every $F\cap E(G[U_i])$ is a forest.
For $d\ge2$,
\[
 \Pr(\mathcal F)
 \le\frac12+\exp\!\left(-\frac{k}{2^d(2^{(d-1)(2^d-1)}-1)}\right).
\]
For $d\le1$, one has $\Pr(\mathcal F)\le2^{-k}$.
\end{corollary}

\begin{proof}
Contract each connected induced region to a marked vertex and delete the resulting internal loops.
Retain all other vertices and edges, obtaining a connected quotient $Q$.
Projection maps $\mathcal Z(G)$ onto $\mathcal Z(Q)$.
Its kernel is the direct sum of the cycle spaces of the induced regions.
Thus the quotient pattern is uniform.
Conditional on that pattern, the internal choices in different regions are independent affine cycle-space choices.

If a marked quotient vertex is closed, the corresponding internal choice is a uniform even subgraph of $G[U_i]$.
Its only forest is the empty edge set.
Since that region has positive rank, its forest probability is at most one half.
For a nonclosed marked vertex, bound the conditional forest probability by one.
If $Z$ is the number of closed marked vertices, these observations give
\[
 \Pr(\mathcal F\mid\text{quotient pattern})\le2^{-Z}.
\]
Any remaining quotient loops are at unmarked vertices.
Delete them before applying the loopless closed-star theorem.
Their independent coordinates do not change any marked closure event.
Average and apply Theorem~\ref{thm:exponential-closed-stars}.
\end{proof}

\paragraph{Quantitative scope.}
For each fixed cut bound $d$, the error decreases exponentially with the number of disjoint cyclic regions.
The explicit constant is large: its binary logarithm is $O(d2^d)$.
This theorem improves the previous inverse-linear dependence on the region count.
It does not by itself establish a polynomial error for unrestricted cycle counts.
When applying it with a growing cut bound, the dependence on $d$ must be retained.
